\section{Related Work}
\label{sec:related}

This section reviews the literature on the components of the pipeline
in Figure~\ref{fig:architecture} and on the modelling problem this
article addresses. We first cover volumetric video delivery and the
transports proposed for it, then subjective quality assessment of
point-cloud video, and finally \gls{qoe} models for volumetric
streaming. Table~\ref{tab:positioning} positions \crystal against the
closest studies.

\begin{table}[!t]
  \caption{Positioning of \crystal against related work.
  \yes{} = supported, \no{} = not addressed, \partially{} = partial.
  ``Network impairment'': delay and loss as controlled independent
  variables (\partially{} denotes bandwidth-only adaptation or offline
  loss injection); ``Reliable \quic'': unmodified, fully reliable \quic
  (\partially{} denotes reliable \gls{tcp} or a modified \quic);
  ``Mechanistic model'': the explanatory variable is derived from
  transport mechanics rather than fitted to delay and loss directly.}
  \label{tab:positioning}
  \centering
  \footnotesize
  \setlength{\tabcolsep}{3pt}
  \renewcommand{\arraystretch}{1.15}
  \begin{tabular*}{\textwidth}{@{\extracolsep{\fill}}llcccccc@{}}
    \toprule
    \thead[l]{Work} & \thead[l]{Codec} &
    \thead{Subjective\\Study} & \thead{VR\\Headset} &
    \thead{Network\\Impairment} & \thead{QoE\\Model} &
    \thead{Reliable\\\quic} & \thead{Mechanistic\\Model} \\
    \midrule
    van der Hooft et al.~\cite{vanderhooft2020qomex} & \vpcc  & \yes & \no        & \partially & \no  & \no        & \no \\
    Subramanyam et al.~\cite{subramanyam2022} & CWI-PCL    & \yes & \yes       & \partially & \no  & \no        & \no \\
    Cox et al.~\cite{Cox2023}               & \vpcc        & \yes & \no        & \no        & \no  & \no        & \no \\
    Nguyen et al.~\cite{nguyen2023}         & \vpcc        & \yes & \partially & \no        & \no  & \no        & \no \\
    Zhou et al.~\cite{zhou2025tomm}         & \vpcc, \gpcc & \yes & \yes       & \no        & \no  & \no        & \no \\
    Van Damme et al.~\cite{vandamme2024}    & \vpcc        & \yes & \no        & \no        & \yes & \no        & \no \\
    Dumic \& da Silva Cruz~\cite{dumic2023} & \vpcc        & \yes & \no        & \partially & \no  & \no        & \no \\
    Ravuri et al.~\cite{ravuri2023}         & \vpcc        & \no  & \no        & \yes       & \no  & \partially & \no \\
    Rudolph et al.~\cite{rudolph2025rabbit} & \vpcc        & \no  & \no        & \no        & \no  & \no        & \no \\
    Udora et al.~\cite{qoeholographic,udora2025thesis} & uncompressed & \yes & \yes & \yes & \yes & \partially & \no \\
    \mosaic~\cite{mosaic2026}               & \vpcc        & \yes & \yes       & \yes       & \no  & \yes       & \no \\
    \vquest~\cite{vquest2026}               & frame-indep. & \yes & \yes       & \yes       & \no  & \yes       & \no \\
    \midrule
    \textbf{\crystal (ours)}                & \vpcc        & \yes & \yes       & \yes       & \yes & \yes       & \yes \\
    \bottomrule
  \end{tabular*}
\end{table}

\subsection{Volumetric Video Delivery and Transport}
\label{sec:rw-systems}

Surveys of immersive media delivery describe the end-to-end
pipeline from capture to \gls{6dof} rendering and identify \gls{qoe}
under network impairment as an open
problem~\cite{Immersivevideo,Viola2023,shen2023toward,garcia22}; Liu et
al.~\cite{pcstreaming21} and Enenche et al.~\cite{metaverseroad25} note
that measuring the \gls{qoe} of point-cloud streaming ``is still in its
infancy''. Streaming systems for \vpcc address bandwidth and
adaptation. Rudolph et al.~\cite{rudolph2025rabbit} transcode \vpcc
streams in real time by re-encoding only the video sub-streams; Shi et
al.~\cite{qv424} add viewpoint-aware tiling to \vpcc using the
encoder's intermediate outputs, with a \gls{qoe} model trained on
VOLVQAD ratings. Codec-agnostic systems exploit inter-frame
correlation~\cite{hermes}, neural super-resolution~\cite{yuzu},
codec-independent progressive partitions~\cite{p2vs},
multiple-description coding over \gls{webrtc}~\cite{mathias24} and
multi-codec rate adaptation~\cite{codec}. These works evaluate with
objective metrics or simulation and hold the encoder configuration, in
particular the \gof length, fixed.

On the transport side, Ravuri et al.~\cite{ravuri2023} integrate
\gls{dash} with partially reliable \quic for \vpcc, deciding per object
whether to deliver reliably, and report 46\% higher throughput and 33\%
fewer playout interruptions than reliable \quic at 5\% loss on 5G
traces, with a headless player. Haseeb et al.~\cite{haseeb2025} combine
adaptive encoding with a mixed-reliability \quic protocol, Freeman et
al.~\cite{moq2025} obtain implicit server-side adaptation from
Media-over-\quic delivery timeouts, Thomas and
Xylomenos~\cite{thomas2025} stream point clouds with forward error
correction on a real 5G testbed without \quic, and Fiorina et
al.~\cite{fiorina2025} estimate \gls{qoe} for multicast point-cloud
streaming over 5G with an objective model. Each either modifies the
fully reliable protocol that production stacks ship~\cite{rfc9000} or
omits subjective evaluation; the perceptual cost of reliable \quic
turning loss into delay has not been quantified. Before~\cite{mosaic2026},
no study combined real \vpcc transmission over \quic, playback in a
\gls{vr} headset, controlled impairment on a real link, 5G delivery and
a standardised subjective study, and no study yet combines those with a
model.

\subsection{Subjective Quality Assessment of Point-Cloud Video}
\label{sec:rw-subjective}

Subjective studies of point-cloud video have concentrated on
compression and bandwidth. Van der Hooft et al.~\cite{vanderhooft2020qomex}
conduct the first subjective evaluation of adaptive \vpcc streaming,
with 30 participants rating pre-rendered viewports on a flat screen
under bandwidth and buffer variation at a fixed latency. Subramanyam
et al.~\cite{subramanyam2022} evaluate tiled user-adaptive streaming
for \gls{vr} remote communication with 33 users in head-mounted
displays, varying bandwidth. Cox et al.~\cite{Cox2023} release VOLVQAD,
7,680 ratings of 376 \vpcc sequences from 120 participants on 2D
displays, and find texture compression more damaging than geometry.
Nguyen et al.~\cite{nguyen2023} characterise \gls{qoe} and immersion of
\vpcc point clouds on the HoloLens~2 with 32 participants across
quantisation, quality switches, viewing distance and content. Zhou et
al.~\cite{zhou2025tomm} collect quality scores and eye-tracking maps
from 40 participants viewing \vpcc, \gpcc and CWI-PCL content with
\gls{6dof} in \gls{vr}. Alexiou et al.~\cite{Alexiou2023} review the
methodology; Ak et al.~\cite{ak2024basics} benchmark objective metrics
on a broad static corpus, where \pcqm~\cite{pcqm} is among the most
stable across codecs, and Liu et al.~\cite{liu2021tip} fit a
reduced-reference perceptual model to \vpcc quantisation parameters.

Few studies vary the network. Dumic and da Silva Cruz~\cite{dumic2023}
inject \SIrange{0.5}{2}{\percent} packet loss into \vpcc sub-bitstreams
before decoding and report \mos drops exceeding one scale unit at
\SI{0.5}{\percent}, with geometry and occupancy losses far more
damaging than attribute losses; Wu et al.~\cite{wu2021nossdav} study
distortion and concealment under loss. In both, loss is applied to the
bitstream rather than produced by a transport, and the content is
shown as 2D video. Our earlier study~\cite{mosaic2026} rated \vpcc in
head-mounted displays under real \quic impairment but paired each
network tier with a different encoder configuration, so the two
factors could not be separated. The only subjective study under delay
and loss on reliable transport in head-mounted displays is otherwise
that of Udora et al., discussed next.

\subsection{\gls{qoe} Models for Volumetric Streaming}
\label{sec:rw-models}

Parametric \gls{qoe} models relate a network impairment to a quality
score. The exponential form used in \crystal descends from the IQX
hypothesis of Fiedler et al.~\cite{fiedler2010}, which derives
$\mathrm{QoE} = a\,e^{-bx} + c$ for a single impairment $x$; Reichl et
al.~\cite{reichl2010} argue for a logarithmic law, and the two agree in
the operating region of interest~\cite{alreshoodi2013}. For volumetric
media, three kinds of model have been proposed. Streaming optimisers
such as QV4~\cite{qv424} and Wang et al.~\cite{wang2022tmm} embed a
weighted sum of quality, quality fluctuation and rebuffering as their
objective; these are calibrated on 2D-video conventions and evaluated
in simulation. Feature-based models map content or device features to
quality without a network variable: Van Damme et
al.~\cite{vandamme2022,vandamme2024} train a clustering-based
psychometric model on no-reference image features for \vpcc video at
video, \gof and frame granularity, and Zhang et al.~\cite{zhang2019towards}
build a fuzzy inference system for the HoloLens.

The third kind, and the closest to this work, is the network-aware
regression of Udora et al.~\cite{qoeholographic,udora2025thesis}. From a
double-stimulus study with 25 participants in head-mounted displays
they fit
\begin{equation}
  \mathrm{QoE} = a\,e^{b D} + c\,\plr + d,
  \label{eq:udora}
\end{equation}
with the coefficients re-estimated for each of three point-cloud
density levels. The transport is reliable (\gls{tcp}), so the data
carry the same loss-to-latency signature as ours; the form is
additive in delay and loss; and held-out \gls{plcc} is 0.57--0.66
across quality levels. \crystal keeps the exponential kernel but
replaces the five fitted coefficients with one perceptual parameter
and a measured ceiling, derives its explanatory variable from
transport mechanics rather than fitting delay and loss directly, and
treats separability as a prediction to be tested.
